\section{Related Work}
\label{sec:related}

\subsection{Prediction Markets and Market Efficiency}

Prediction markets have a rich history as forecasting instruments, with theoretical foundations in the Efficient Market Hypothesis (EMH) and practical applications ranging from election forecasting to corporate decision-making~\cite{arrow2008promise}. Wolfers and Zitzewitz~\cite{wolfers2004prediction} provide a comprehensive analysis of prediction market accuracy, demonstrating that market prices approximate true probabilities under reasonable conditions. Manski~\cite{manski2006interpreting} offers a more nuanced view, showing that market prices can deviate from mean beliefs when risk preferences are heterogeneous.

Modern decentralized prediction markets such as Polymarket operate on blockchain infrastructure with automated market maker (AMM) pricing, enabling 24/7 trading without centralized counterparties~\cite{chen2007utility}. These markets exhibit several well-documented inefficiencies: the favorite-longshot bias, whereby low-probability events are systematically overpriced~\cite{snowberg2010explaining}; recency bias, where recent information is overweighted relative to base rates~\cite{rothschild2009forecasting}; and illiquidity effects, where thin markets can sustain substantial mispricings for extended periods. Our Context Planner is designed specifically to exploit the latter two inefficiencies by injecting base-rate information and domain-specific context that thin markets may not have incorporated.

\subsection{LLM-Based Forecasting}

Recent work has explored the use of large language models as forecasting tools. Halawi et al.~\cite{halawi2024approaching} demonstrate that retrieval-augmented LLMs can approach human-level forecasting accuracy on Metaculus questions, achieving calibration comparable to competitive human forecasters when provided with relevant context. Schoenegger et al.~\cite{schoenegger2024wisdom} compare LLM forecasts against human crowds, finding that while LLMs exhibit reasonable calibration on well-known event types, they struggle with novel or rapidly evolving situations.

A critical limitation of existing LLM forecasting approaches is the absence of \emph{domain-specific context injection}. Generic retrieval-augmented generation retrieves broadly relevant documents but does not identify what specific information the LLM \emph{lacks} for a given prediction domain. For commodity price predictions, the LLM needs current spot prices and inventory data; for award predictions, it needs precursor ceremony results; for geopolitical events, it needs recent diplomatic communications. Our Context Planner addresses this gap through domain classification and targeted knowledge gap identification, a capability absent from prior work.

\subsection{Agent-Based Market Simulation}

Agent-based computational economics has a substantial history, originating with the Santa Fe Artificial Stock Market~\cite{arthur1997asset} and the zero-intelligence trader framework of Gode and Sunder~\cite{gode1993allocative}. These models demonstrate that market-level phenomena such as price discovery, volatility clustering, and bubble formation can emerge from simple agent interactions without requiring individual rationality.

More recently, LLM-powered agents have been applied to financial simulation. FinAgent~\cite{zhang2024finagent} uses LLMs for trading decision-making in equity markets, while TradingGPT~\cite{li2023tradinggpt} explores multi-agent financial interaction. However, these systems focus on \emph{using} agents for trading rather than \emph{discovering} trading strategies through simulation. Our Virtual Prediction Market fills this gap by using heterogeneous agent populations as a testbed for strategy discovery, where candidate strategies must survive adversarial interaction with noise traders, momentum followers, and contrarian agents.

\subsection{Automated Scientific Discovery and Code Evolution}

The concept of automating scientific research has gained momentum through several recent advances. FunSearch~\cite{romera2024mathematical} demonstrates that LLM-driven code evolution can discover novel mathematical constructions, using an evolutionary loop where LLMs propose program modifications evaluated against formal criteria. The autoresearch paradigm, articulated by Karpathy and instantiated in systems such as AI Scientist~\cite{lu2024aiscientist}, extends this approach to experimental design and hypothesis generation.

Our Evolution Engine applies these principles to prediction strategy discovery. Unlike parameter optimization approaches such as Bayesian optimization or grid search, evolutionary code search can discover qualitatively new analysis patterns---for instance, learning to weight information sources differently based on market liquidity or to apply asymmetric adjustments depending on whether the market price falls in overconfident or underconfident ranges. This capability is essential for prediction markets, where the optimal strategy varies across domains and evolves as market participants adapt.
